\section*{Chapitre \ref{ch:rc:margin-models} --- Modèles de marge}

\begin{solution}{exo:rc:margin-models:1}
$4~\%\times100 = 4$ millions pour l'échéance de dix ans et $1~\%\times100 = 1$ million pour celle d'un an~: 5~millions de dollars bruts.
\end{solution}

\begin{solution}{exo:rc:margin-models:2}
Son quantile à 99~\% est le troisième plus mauvais parmi environ 250 mouvements de cinq jours~: quelques jours volatils récents ajoutent quelques
nouvelles observations de queue sans modifier beaucoup le quantile, tant qu'elles ne dépassent pas en nombre l'ancienne queue, et les jours calmes
de l'année passée dominent encore la fenêtre.
\end{solution}

\begin{solution}{exo:rc:margin-models:3}
$\mathrm{NGR} = 30/120 = 0.25$~; marge nette $= (0.4+0.6\times0.25)\times\text{brut} = 0.55\times\text{brut}$~: une réduction de 45~\%.
\end{solution}

\begin{solution}{exo:rc:margin-models:4}
Les pondérations de risque du SIMM sont calibrées sur de longues séries historiques incluant des périodes de tension, et elles ne changent
qu'à la recalibration~; elles ne réagissent pas à la volatilité courante. Les modèles des chambres de compensation suivent la volatilité (filtrés
ou à courte fenêtre d'observation), si bien que leur marge a monté avec elle.
\end{solution}

\begin{solution}{exo:rc:margin-models:5}
Le pic est fixé par la marge du modèle au plus fort de la tension, que les outils ne modifient pas (le
coussin est libéré, la pondération sous tension étant inférieure à son niveau de mars)~; ils ne font qu'augmenter la marge détenue
auparavant, de sorte que la hausse est plus faible.
\end{solution}

\begin{solution}{exo:rc:margin-models:6}
L'impact en temps et en prix de la liquidation d'une position importante par rapport au volume quotidien~: la clôture
prend plus de temps que la période de marge en risque et déplace les prix, ce qu'un modèle de mouvements de taille normale sur un
horizon fixe ne voit pas.
\end{solution}

\begin{solution}{exo:rc:margin-models:7}
Exceptions attendues~: $2\,177\times1~\% = 21.8$~; rapport de vraisemblance de Kupiec 6.86, p-valeur 0.009. Des fenêtres de cinq jours qui se chevauchent font que
les jours consécutifs partagent quatre de leurs cinq jours de mouvements, de sorte que les exceptions arrivent par séries~; le test suppose
l'indépendance et surestime la preuve.
\end{solution}

\begin{solution}{exo:rc:margin-models:8}
La couverture doit être testée hors échantillon et en période de tension~: un modèle calibré sur des mois calmes couvre des mois calmes
et échoue quand la volatilité bondit, comme en mars 2020. Le test porte sur le mois suivant, et non sur le dernier.
\end{solution}

\begin{solution}{pb:rc:margin-models:1}
\textbf{1.} 4.21~millions de dollars en février~; 10.30~millions de dollars le 23 mars.
\textbf{2.} 6.09~millions de dollars.
\textbf{3.} Volatilité~: les positions sont figées dans le rejeu.
\textbf{4.} 1.06~million de dollars (de 4.19 à 5.25).
\textbf{5.} Elle couvre le risque courant~: le modèle historique a sous-estimé la marge du portefeuille pendant la tension,
qui est précisément le moment où un défaut est le plus probable.
\textbf{6.} 5.26~millions de dollars de marge détenue~; une hausse de 5.04~millions de dollars.
\textbf{7.} 4.79~millions de dollars de marge détenue~; une hausse de 5.51~millions de dollars.
\textbf{8.} Le financement de la marge supplémentaire détenue les années calmes (environ 1~million sous le coussin, 0.6 sous
la pondération en tension de février), au spread de financement du membre au-dessus de ce que rapporte la marge.
\textbf{9.} Après le retour de la volatilité à la normale, et progressivement~: la reconstituer d'un coup serait un nouvel
appel de marge au moment même où les marchés se redressent.
\textbf{10.} D'une hausse de 25~\% sur cinq jours à une hausse pouvant atteindre 80~\%.
\textbf{11.} Historique sur dix jours~: 4.15~millions de dollars~; sensibilités~: 8.98~millions de dollars~; barème~: 13.19~millions de dollars
brut, 9.23~millions de dollars net.
\textbf{12.} Netting sur l'ensemble de ses transactions compensées, marge plus faible que le barème, et aucun risque bilatéral
de contrepartie~; en échange, une exposition au modèle procyclique de la chambre de compensation.
\textbf{13.} Tous les montants sont inférieurs à 50~millions d'euros~: sous le seuil, aucune marge initiale ne serait
échangée bilatéralement pour ce seul portefeuille.
\textbf{14.} En crise, les deux augmentent~: la marge de variation suit les cours au jour le jour et doit être versée en espèces~;
la marge initiale monte avec la volatilité. Ensemble, elles assèchent la liquidité au pire moment.
\textbf{15.} Le modèle propre de la chambre de compensation (ou une réplique) exécuté sur des scénarios de tension, et non la marge des périodes calmes.
\textbf{16.} Un atout pour la sécurité de la chambre de compensation, un défaut pour la liquidité du système~: elle transfère
la tension aux membres au moment où ils peuvent le moins la supporter.
\textbf{17.} Les membres, par des marges plus élevées en période calme, qu'ils doivent financer.
\textbf{18.} Par la hausse de la marge dans les scénarios de tension historiques et hypothétiques sévères (ici plus de 6~millions de dollars
en trois semaines), plus la marge de variation, avec un collatéral livrable le jour même.
\textbf{19.} Résultat nommé~: \emph{les appels de marge de mars 2020}~: le modèle filtré a appelé 6.09~millions de dollars
de marge nouvelle~; avec un coussin de 25~\%, 5.04~millions de dollars, et avec une pondération en tension de 25~\%, 5.51~millions de dollars.
\textbf{20.} Protection de la contrepartie contre la prévisibilité et la liquidité du payeur.
\end{solution}

\begin{solution}{iq:rc:margin-models:1}
Elle couvre la perte de clôture du portefeuille d'un défaillant sur la période de marge en risque, à un niveau de
confiance élevé. Les chambres de compensation utilisent des modèles historiques, filtrés ou fondés sur des scénarios, appliqués à l'ensemble du portefeuille du membre,
au moins 99~\% sur quelques jours.
\iqlookfor{objet, horizon, niveau de confiance et méthode.}
\end{solution}

\begin{solution}{iq:rc:margin-models:2}
La marge monte avec la volatilité, si bien que les membres doivent déposer davantage précisément en période de tension, ce qui force des ventes et
alimente la tension. Outils~: coussins libérés en tension, pondérations en période de tension, longues fenêtres d'observation et
planchers de volatilité.
\iqlookfor{la boucle de rétroaction et les trois outils.}
\end{solution}

\begin{solution}{iq:rc:margin-models:3}
On calcule les sensibilités par \omterm{def:rc:market-risk-measures:factor}{facteur de risque} (delta par échéance, vega, convexité), on les multiplie par les pondérations de risque calibrées,
on agrège au sein des compartiments avec des corrélations, puis entre compartiments et classes~; les paramètres sont
recalibrés périodiquement sur des historiques incluant des périodes de tension.
\iqlookfor{sensibilités, pondérations, corrélations et calibration.}
\end{solution}

\begin{solution}{iq:rc:margin-models:4}
Des modèles différents (historiques, filtrés, fondés sur des scénarios), des fenêtres d'observation, des périodes de marge en risque, des niveaux de confiance,
des \omterm{def:rc:margin-models:apc}{outils anti-procyclicité}, des \omterm{def:rc:margin-models:conc}{majorations pour concentration} et pour liquidité, et des \omterm{def:rc:counterparty-exposure:netting}{ensembles de compensation}.
\iqlookfor{différences de modèles et de paramètres.}
\end{solution}

\begin{solution}{iq:rc:margin-models:5}
On compare les pertes réalisées sur la période de marge en risque à la marge détenue, par portefeuille, et on teste
la fréquence et le regroupement des exceptions~; c'est difficile, car les portefeuilles changent, les fenêtres se chevauchent, les tensions sont rares et la
perte pertinente inclut les coûts de liquidation.
\iqlookfor{la procédure et ses difficultés.}
\end{solution}

\begin{solution}{iq:rc:margin-models:6}
Les appels de marge, initiale et de variation, ont été importants, rapides et synchronisés~; les établissements qui ne les avaient pas anticipés
ont vendu des actifs sur des marchés en baisse. Les modèles de marge doivent être prévisibles, et les établissements ont besoin de plans de liquidité
fondés sur la marge en période de tension.
\iqlookfor{effets sur la liquidité et prévisibilité.}
\end{solution}
